% Generated draft: requires mathematical and source review.
\documentclass[11pt]{article}
\usepackage[margin=25mm]{geometry}
\usepackage{fontspec}
\setmainfont{Latin Modern Roman}
\usepackage{amsmath,amssymb,mathtools,unicode-math}
\setmathfont{Latin Modern Math}
\usepackage{enumitem,needspace,xurl,hyperref}
\setlength{\parindent}{0pt}
\setlength{\parskip}{5pt}
\setlength{\emergencystretch}{3em}
\newcommand{\N}{\mathbb{N}}
\newcommand{\Z}{\mathbb{Z}}
\newcommand{\Q}{\mathbb{Q}}
\newcommand{\R}{\mathbb{R}}
\newcommand{\C}{\mathbb{C}}
\newcommand{\F}{\mathbb{F}}
\newcommand{\A}{\mathbb{A}}
\newcommand{\PP}{\mathbb{P}}
\newcommand{\E}{\mathbb{E}}
\newcommand{\eps}{\varepsilon}
\newcommand{\dd}{\,\mathrm{d}}
\begin{document}
\section*{UnsolvedMath ID: 701196}
\textbf{Source problem code:} catalogue entry
\subsection{Definitions and mathematical statement}
For a finite group \(G\) with identity \(1_G\), the conjugacy class of \(a\in G\) is
\[
a^G=\{g^{-1}ag:g\in G\}.
\]
For subsets \(C,D\subseteq G\), define their setwise product by
\[
CD=\{cd:c\in C,\ d\in D\}.
\]
A group is simple if its only normal subgroups are \(\{1_G\}\) and \(G\). The question asks whether there exist a finite non-abelian simple group \(G\) and elements \(a,b\in G\setminus\{1_G\}\) such that
\[
a^G b^G=(ab)^G.
\]
Equivalently, can two nontrivial conjugacy classes \(C,D\) of a finite non-abelian simple group have a product that is a single conjugacy class? The classes \(C,D\) are not required to be distinct. The source does not define ``nontrivial'': it may mean different from \(\{1_G\}\), or having more than one element. These interpretations coincide for non-abelian simple groups, whose center is trivial, so the displayed existence question covers either interpretation of the intended target.
\subsection{Short English statement}
Can a finite non-abelian simple group have two nontrivial conjugacy classes whose setwise product is one conjugacy class? This is attributed to Z. Arad, Kourovka Notebook Problem 16.4 (2006). The supplied registry labels it open; that label has not been independently verified.
\subsection{Sources}
\begin{itemize}
\item UnsolvedMath local registry, permanent ID \(701196\), exact source problem code \texttt{catalogue entry}, title \emph{Kourovka Notebook Problem 16.4}. The supplied record attributes the problem to Z. Arad, with proposed year 2006, and labels its status \texttt{open}.
\item Supplied source attribution: \emph{The Kourovka Notebook (Unsolved Problems in Group Theory)}, 16th issue (2006), Problem 16.4; also identified in the record with the 21st issue edition (2026), \texttt{arXiv:1401.0300v48}. These bibliographic details are reproduced from the supplied record and have not been independently checked.
\item The supplied collection description identifies E. I. Khukhro and V. D. Mazurov as editors of the Kourovka Notebook collection covering Issues 1--21 (1965--2026).
\end{itemize}
\end{document}
